\documentclass[11pt,a4paper]{article}

% === 页面与字体 ===
\usepackage[margin=2.5cm]{geometry}
\usepackage{fontspec}
\setmainfont{TeX Gyre Termes}
\usepackage{xeCJK}
\setCJKmainfont{Noto Serif SC}
\setCJKfamilyfont{song}{Noto Serif SC}
\newcommand{\hei}{\CJKfamily{song}\bfseries}

% === 表格 ===
\usepackage{booktabs}
\usepackage{tabularx}
\usepackage{multirow}
\usepackage{array}
\usepackage{dcolumn}
\newcolumntype{d}[1]{D{.}{.}{#1}}

% === caption 中文格式 ===
\usepackage{caption}
\renewcommand{\tablename}{表}
\captionsetup[table]{labelsep=quad, font={bf}, justification=centering}

% === 其他 ===
\usepackage{amsmath}
\usepackage{hyperref}
\hypersetup{colorlinks=false}
\usepackage{setspace}
\onehalfspacing
\usepackage{enumitem}

\usepackage{fancyhdr}
\pagestyle{fancy}
\fancyhf{}
\rhead{\small 数据要素利用与资本市场定价效率}
\lhead{\small 内生性检验}
\cfoot{\thepage}

\title{{\hei\Large 数据要素利用与资本市场定价效率} \\[8pt]
{\large 内生性检验（v2：统一口径）}}
\author{}
\date{2026年3月}

\begin{document}
\maketitle
\thispagestyle{fancy}

% ============================================================
\section*{一、内生性问题与处理策略}

在基准回归中，数据要素利用程度（$DU\_kw$）对股价延迟（$PriceDelay$）的负向效应在统计和经济意义上均显著（$\beta=-0.0040$, $t=-3.94$）。然而，该估计可能受到遗漏变量偏误和反向因果的困扰。为缓解上述问题，本文采用以下策略：（1）同行业均值工具变量法；（2）Bartik移位份额工具变量法；（3）滞后一期自变量；（4）Oster (2019) 系数稳定性检验。所有检验均使用同期$DU\_kw$口径（与基准回归一致），以企业和年份双向固定效应、行业$\times$年份聚类标准误为标准设定。

% ============================================================
\section*{二、工具变量法}

\subsection*{（一）同行业均值工具变量}

参照已有文献的做法（Chang等, 2024; 李世刚等, 2025; 朱康和汤甬, 2025），本文以同年度同行业（剔除本企业）的数据要素利用均值作为工具变量。该工具变量的逻辑在于：同行业企业面临相似的技术环境和竞争压力，行业层面的数据要素利用趋势影响个体企业的数据化进程，但行业均值本身不太可能直接影响某一特定企业的股价延迟，从而满足排他性约束。工具变量$IV_{Peer,i,t}$的构造方式为：
\begin{equation}
IV_{Peer,i,t} = \frac{1}{N_{j,t}-1}\sum_{k \in j, k \neq i} DU\_kw_{k,t}
\end{equation}

两阶段最小二乘法（2SLS）结果如表\ref{tab:iv}列(1)所示。Kleibergen-Paap rk Wald F统计量为796.67，远超Stock和Yogo (2005) 在10\%最大偏误下的临界值16.38，排除弱工具变量的可能。Anderson-Rubin弱工具变量稳健检验的F值为5.23（$p=0.022$），表明即使在弱工具变量情形下，$DU\_kw$对$PriceDelay$的效应仍然显著。第二阶段回归中$DU\_kw$系数为$-0.0104$（$t=-2.31$），在5\%水平上显著为负，方向和显著性与OLS一致，且系数绝对值更大。

此外，Durbin-Wu-Hausman检验的F值接近零（$p=0.988$），不拒绝$DU\_kw$外生的原假设。这意味着OLS估计量本身可能就是一致的，内生性问题并不严重，IV结果为OLS提供了额外的稳健性保障。

\subsection*{（二）Bartik移位份额工具变量}

借鉴Goldsmith-Pinkham等 (2020) 的方法，本文构造Bartik型移位份额工具变量。以基期（2011年）各行业的数据要素利用水平作为"份额"，行业整体增长率作为"移位"，两者交互构成外生预测：
\begin{equation}
IV_{Bartik,i,t} = \overline{DU\_kw}_{j,2011} \times g_{j,t}
\end{equation}
其中$\overline{DU\_kw}_{j,2011}$为行业$j$基期均值，$g_{j,t}$为行业从基期到$t$期的增长率。该工具变量的排他性建立在以下假设之上：基期的行业数据化水平通过行业发展趋势影响当期个体企业的数据化程度，但基期行业特征不直接影响当期个体企业的股价延迟。

表\ref{tab:iv}列(2)显示，Kleibergen-Paap F统计量高达3342.80，工具变量极强。Anderson-Rubin F为5.71（$p=0.017$）。第二阶段$DU\_kw$系数为$-0.0102$（$t=-2.43$），在5\%水平上显著为负，与同行业均值IV结果高度一致。DWH检验同样不拒绝外生性（$p=0.602$）。

% ============================================================
\section*{三、滞后一期自变量}

为进一步缓解反向因果担忧，本文将自变量替换为$DU\_kw_{t-1}$。表\ref{tab:iv}列(3)显示，OLS估计中$DU\_kw_{t-1}$系数为$-0.0036$（$t=-3.39$），在1\%水平上显著为负。以$DU\_kw_{t-1}$的同行业均值滞后值作为工具变量，2SLS估计系数为$-0.0106$（$t=-1.89$），在10\%水平上显著，一阶段F统计量为422.92。滞后一期的结果与同期口径方向完全一致，但样本量因丢失首期观测而减少（$N=40,326$）。

% ============================================================
\section*{四、Oster系数稳定性检验}

Oster (2019) 提出基于系数稳定性和$R^2$变化评估遗漏变量偏误的方法。该方法的核心思想是：如果加入控制变量后系数变化很小而$R^2$提升明显，则遗漏变量不太可能推翻结论。

本文计算结果如下：不含控制变量时，$\hat{\beta}=-0.0060$，$R^2=0.408$；加入全部控制变量后，$\hat{\beta}=-0.0040$，$R^2=0.432$。以$R^2_{max}=\min(1, 1.3\tilde{R})=0.562$计算，$\delta^*=10.7$，即遗漏变量需要达到可观测控制变量的10.7倍才能将系数推至零。以更保守的$R^2_{max}=2\tilde{R}-\dot{R}=0.457$计算，$\delta^*=2.0$，偏误调整后$\hat{\beta}_{adj}=-0.0020$仍为负值。两种设定下$\delta^*$均大于1，满足Oster (2019) 提出的稳健性标准。

% ============================================================
\section*{五、内生性检验结果汇总}

\begin{table}[htbp]
\centering
\caption{内生性检验结果汇总}
\label{tab:iv}
\vspace{6pt}
\begin{tabular}{lcccccc}
\toprule
& (0) & (1) & (2) & (3) & (4) \\
& OLS & Peer IV & Bartik IV & OLS($t{-}1$) & IV($t{-}1$)\\
\midrule
$DU\_kw$ / $DU\_kw_{t-1}$ & $-$0.0040$^{***}$ & $-$0.0104$^{**}$ & $-$0.0102$^{**}$ & $-$0.0036$^{***}$ & $-$0.0106$^{*}$\\
& (0.0010) & (0.0045) & (0.0042) & (0.0011) & (0.0056) \\[6pt]
\textit{Panel B: IV诊断} \\[3pt]
KP rk Wald $F$ & & 796.67 & 3342.80 & & 422.92 \\
Anderson-Rubin $F$ & & 5.23$^{**}$ & 5.71$^{**}$ & & \\
DWH内生性检验 $F$ & & 0.00 & 0.27 & & \\
DWH $p$值 & & 0.988 & 0.602 & & \\[3pt]
\textit{Panel C: Oster bounds} \\[3pt]
$\delta^*$ ($R^2_{max}=1.3\tilde{R}$) & \multicolumn{5}{c}{10.7} \\
$\delta^*$ ($R^2_{max}=2\tilde{R}-\dot{R}$) & \multicolumn{5}{c}{2.0} \\
$\hat{\beta}_{adj}$ ($\delta=1$, 保守) & \multicolumn{5}{c}{$-$0.0020} \\[3pt]
\midrule
控制变量 & 是 & 是 & 是 & 是 & 是 \\
企业固定效应 & 是 & 是 & 是 & 是 & 是 \\
年份固定效应 & 是 & 是 & 是 & 是 & 是 \\
$N$ & 45,296 & 45,296 & 45,296 & 40,326 & 40,326 \\
\bottomrule
\end{tabular}

\vspace{6pt}
{\small 注：括号内为行业$\times$年份聚类稳健标准误。$^{***}$、$^{**}$、$^{*}$分别表示在1\%、5\%、10\%水平上显著。KP $F$为Kleibergen-Paap rk Wald $F$统计量。Anderson-Rubin检验为弱工具变量稳健推断。DWH为Durbin-Wu-Hausman内生性检验，原假设为$DU\_kw$外生。Oster $\delta^*$为将系数推至零所需的相对遗漏变量选择度，$\delta^*>1$表示结果稳健。所有回归均使用同期$DU\_kw$口径，列(3)(4)使用$DU\_kw_{t-1}$。}
\end{table}

表\ref{tab:iv}汇总了全部内生性检验结果。可以得到以下结论：第一，所有方法均得到$DU\_kw$对$PriceDelay$的显著负向效应，系数方向完全一致，说明结果不受特定估计方法驱动。第二，两种工具变量的系数绝对值（0.010左右）均大于OLS（0.004），这一模式与测量误差导致的"衰减偏误"预期相符：$DU\_kw$作为关键词计数的代理变量，含有噪声，OLS系数被测量误差向零拉拽，IV修正后效应更大。第三，DWH检验均不拒绝$DU\_kw$的外生性，说明遗漏变量偏误和反向因果并非本研究的主要威胁，IV结果为OLS提供了额外的稳健性保障。第四，Oster $\delta^*$在两种$R^2_{max}$设定下均大于1，遗漏变量需要达到可观测控制变量的2至10.7倍才能推翻结论。综上，内生性检验结果支持数据要素利用对资本市场定价效率具有促进作用这一核心发现的稳健性。

% ============================================================
\section*{参考文献}

\begin{enumerate}[label={[\arabic*]}, leftmargin=2em, itemsep=3pt]

\item Chang Y, Liu C, Zhou F. Digital Transformation and Capital Market Pricing Efficiency: Identification from a Natural Experiment[J]. \textit{Finance Research Letters}, 2024, 65: 105557.

\item Goldsmith-Pinkham P, Sorkin I, Swift H. Bartik Instruments: What, When, Why, and How[J]. \textit{American Economic Review}, 2020, 110(8): 2586--2624.

\item Oster E. Unobservable Selection and Coefficient Stability: Theory and Evidence[J]. \textit{Journal of Business \& Economic Statistics}, 2019, 37(2): 187--204.

\item Stock J H, Yogo M. Testing for Weak Instruments in Linear IV Regression[C]// Andrews D W K, Stock J H. Identification and Inference for Econometric Models. Cambridge University Press, 2005: 80--108.

\item 李世刚, 唐松, 施海晶. 企业数据资产信息披露与资本市场定价效率[J]. 中国工业经济, 2025(1): 137--155.

\item 朱康, 汤甬. 数据要素利用与企业金融资产配置[J]. 会计研究, 2025(2).

\end{enumerate}

\end{document}
